\documentclass{article}

% NeurIPS 2026 workshop style — sty file assumed available in build environment
\usepackage{neurips_2026}

% Standard packages
\usepackage[utf8]{inputenc}
\usepackage[T1]{fontenc}
\usepackage{hyperref}
\usepackage{url}
\usepackage{booktabs}
\usepackage{amsfonts}
\usepackage{amsmath}
\usepackage{amssymb}
\usepackage{nicefrac}
\usepackage{microtype}
\usepackage{xcolor}
\usepackage{graphicx}
\usepackage{listings}
\usepackage{algorithm}
\usepackage{algorithmic}

% Bibliography
\usepackage[numbers]{natbib}

% Convenience macros
\newcommand{\gore}{\textsc{Gore}}
\newcommand{\todo}[1]{\textcolor{red}{\textbf{[TODO: #1]}}}

% -----------------------------------------------------------------------

\title{GORE: A Controlled Environment for Training and Evaluating
       Backtracking Reasoning in Language Models}

\author{%
  % TODO: fill in author block
  Anonymous Author(s) \\
  \texttt{anonymous@example.com}
}

\begin{document}

\maketitle

\begin{abstract}
% TODO: write abstract (~150 words).
% Key points to hit:
%   - LLMs struggle with multi-step backtracking / systematic search
%   - GORE: minimal Turing-complete graph-rewriting language with explicit backtracking semantics
%   - Fully synthetic, curriculum-controlled training data
%   - SFT baseline + GRPO with dense reward signal
%   - Spectral-R1: representation-level analysis of the learned reasoning process
%   - Transfer signal to natural-language reasoning benchmarks
\todo{Write abstract}
\end{abstract}

% -----------------------------------------------------------------------
\section{Introduction}
\label{sec:intro}

% TODO: motivate the problem of backtracking reasoning in LLMs.
% - Transformers as "forward-only" reasoners; failure modes on search/planning tasks
% - Why controlled synthetic environments matter (cf. BabyAI, ARC, etc.)
% - GORE as a minimal Turing-complete substrate: designed so that *every* non-trivial
%   program requires backtracking to evaluate
% - Overview of contributions:
%     1. The GORE language and interpreter
%     2. Synthetic curriculum generator
%     3. Training experiments (SFT, GRPO)
%     4. Spectral-R1 analysis framework
%     5. NL transfer results
\todo{Write introduction}

% -----------------------------------------------------------------------
\section{The \gore{} Language}
\label{sec:language}

\subsection{Design Principles}
\label{sec:design}

% TODO: articulate design goals
% - Minimal: fewest primitives needed for Turing completeness
% - Transparent: execution trace is a human-readable sequence of rewrite steps
% - Backtrack-forcing: no program of interest is solvable by linear scan
% - Formally specified: full small-step operational semantics
\todo{Describe design principles}

\subsection{Syntax}
\label{sec:syntax}

% TODO: give BNF / grammar table for GORE
% Reference the GORE repo (https://github.com/Fr0do/gore) for interpreter source
\todo{Define syntax}

\subsection{Operational Semantics}
\label{sec:semantics}

% TODO: small-step semantics as inference rules (mathpar or proof trees)
% Highlight the rules that introduce / resolve choice points (backtrack edges)
\todo{Define operational semantics}

\subsection{Complexity and Expressiveness}
\label{sec:expressiveness}

% TODO: sketch Turing-completeness argument; note relevant complexity classes
% for the benchmark suite (e.g. fixed-depth search is in PSPACE)
\todo{Discuss expressiveness}

% -----------------------------------------------------------------------
\section{Synthetic Data Generation}
\label{sec:data}

\subsection{Program Sampler}
\label{sec:sampler}

% TODO: describe the parameterized generator
% - Knobs: graph depth D, branching factor B, backtrack frequency F
% - Rejection sampling to hit target difficulty buckets
\todo{Describe program sampler}

\subsection{Curriculum}
\label{sec:curriculum}

% TODO: curriculum schedule: easy → hard over training
% - Stage definitions (e.g., D=2,B=2 → D=4,B=4 → D=6,B=6)
% - How difficulty is measured (number of backtrack steps in gold trace)
\todo{Describe curriculum}

\subsection{Trace Format}
\label{sec:trace}

% TODO: describe the token format fed to the model
% - <think> block containing step-by-step rewrite trace
% - Final answer token
% - Example (small program, full trace)
\todo{Describe trace format and show an example}

% -----------------------------------------------------------------------
\section{Experiments}
\label{sec:experiments}

\subsection{E1: Supervised Fine-Tuning Baseline}
\label{sec:e1}

% TODO:
% - Model(s): base LM (e.g. Qwen-2.5-7B or similar)
% - Training setup: dataset size, curriculum schedule, hyperparameters
% - Metrics: exact-match on held-out GORE programs; backtrack accuracy
% - Ablations: with/without curriculum, trace length analysis
\todo{Report SFT baseline results}

\subsection{E2: GRPO with Dense Backtracking Reward}
\label{sec:e2}

% TODO:
% - Reward design: partial credit per correct rewrite step (dense),
%   +1 for correct final answer, 0 otherwise
% - GRPO training loop, KL budget, group size G
% - Comparison to E1; learning curves
% - Emergent backtracking behavior (qualitative examples)
\todo{Report GRPO results}

\subsection{E3: Spectral-R1 Analysis}
\label{sec:e3}

% TODO:
% - Spectral-R1 framework: PCA / SVD of residual-stream activations
%   conditioned on reasoning state (forward pass vs. backtrack step)
% - Do backtrack steps occupy a distinguishable subspace?
% - Intervention experiments: ablate the "backtrack direction" → performance drop?
% - Connection to broader Spectral-R1 research agenda
\todo{Report spectral analysis results}

\subsection{E4: Transfer to Natural-Language Reasoning}
\label{sec:e4}

% TODO:
% - Evaluate GORE-trained model (or adapter) on NL benchmarks
%   (e.g., GSM8K, MATH, ARC-Challenge, or a custom backtracking-heavy subset)
% - Zero-shot and few-shot settings
% - Compare to equivalent compute on NL data only
\todo{Report NL transfer results}

% -----------------------------------------------------------------------
\section{Related Work}
\label{sec:related}

% TODO: cover the following clusters (1–2 paragraphs each):
%
% 1. Synthetic reasoning environments
%    - BabyAI, ARC, CLUTRR, Countdown, etc.
%
% 2. Chain-of-thought and scratchpad reasoning
%    - Wei et al. (2022), Nye et al. (2021), Lightman et al. (2023)
%
% 3. Reinforcement learning for LLM reasoning
%    - RLHF (Ouyang et al.), GRPO (Shao et al.), DeepSeek-R1, DAPO
%
% 4. Backtracking / search in language models
%    - Tree-of-Thoughts (Yao et al.), RAP, Monte-Carlo Tree Search for LLMs
%
% 5. Mechanistic interpretability of reasoning
%    - Logit lens, probing, linear representation hypothesis
%    - Spectral / SVD analyses of transformer internals
\todo{Write related work}

% -----------------------------------------------------------------------
\section{Conclusion}
\label{sec:conclusion}

% TODO:
% - Summary of contributions
% - Limitations: GORE is synthetic; transfer gap may exist for real-world tasks
% - Future work: larger models, richer reward shaping, multi-agent GORE
\todo{Write conclusion}

% -----------------------------------------------------------------------
\section*{Acknowledgements}
% TODO: add acknowledgements / funding statement before camera-ready
\todo{Add acknowledgements}

% -----------------------------------------------------------------------
\bibliographystyle{plainnat}
\bibliography{references}

\end{document}
